\documentclass{article}
\usepackage{iclr2027_conference,times}
\usepackage[utf8]{inputenc}
\usepackage[T1]{fontenc}
\usepackage{hyperref,url,booktabs,longtable,amsmath,amssymb}
% Draft-only labels; remove these patches for an actual submission.
\usepackage{etoolbox}
\makeatletter
\patchcmd{\@maketitle}{Under review as a conference paper at ICLR 2027}{Working draft for ICLR 2027 --- not submitted}{}{}
\patchcmd{\@maketitle}{Paper under double-blind review}{Author details pending}{}{}
\makeatother
\hypersetup{hidelinks}
\title{REA: Recursive Engineering Agents with Versioned Tool, Parts, and Dataset Memory}
\author{Anonymous authors}
\begin{document}
\maketitle
\begin{abstract}
Engineering agents must combine procedural knowledge with evidence about tools, components, datasets, and their compatibility. This working draft presents Recursive Engineering Agents (REA), an experimental adaptation of procedural scientific skills to recursive language model execution over versioned external memory. REA separates skills, tools, parts, datasets, and supporting evidence, and provides specialist profiles covering 71 engineering domains and 609 subcategories. Explicit category membership scopes catalog access while preserving source references. The current implementation supports bounded recursive execution and immutable record addressing; it does not provide an unlimited model context window. We report implementation checks and a catalog census, and specify a comparative evaluation protocol. Engineering task accuracy, gains over retrieval baselines, and generalization across domains remain unmeasured. This draft is a research scaffold and is not ready for submission.
\end{abstract}

\section{Introduction}
Engineering reasoning depends on more than identifying a plausible procedure. An agent may need to determine whether a solver supports a boundary condition, whether a component meets an interface constraint, or whether a dataset contains the required variables. The relevant evidence may span procedural documents, tool descriptions, component catalogs, and dataset documentation. Providing this material as a single model prompt couples the accessible evidence to the model's finite input capacity.

Scientific Agent Skills organizes procedural knowledge for research agents \citep{Kassis.2026}. Recursive Language Models expose long inputs as an external environment that can be inspected programmatically and processed through recursive model calls \citep{Zhang.2025}. REA combines these mechanisms with versioned engineering catalog records. Its proposed contribution is a reproducible way to scope and trace engineering evidence across recursive investigations. Whether this combination improves engineering outcomes is the central empirical question, not an established result.

The implementation contribution consists of separate memory namespaces, stable record identity, catalog provenance checks, and a complete engineering taxonomy shared between runtime specialist profiles and this appendix. The planned scientific contribution requires demonstrating that these mechanisms improve independently checked task outcomes under matched resource budgets. We do not infer capability from the number of indexed records or the number of specialist profiles.

\section{Related work}
Procedural skill libraries externalize instructions that can be loaded when relevant \citep{Kassis.2026}. Recursive Language Models instead address how a model can inspect and decompose a large external input \citep{Zhang.2025}. REA adopts the latter execution mechanism rather than introducing a new recursive model architecture. Persistent memory systems such as MemGPT address management of information beyond the immediate prompt \citep{Packer.2023}; the distinctive engineering question here concerns versioned evidence, compatibility constraints, and domain scoping.

A scite-assisted discovery pass produced 262 distinct DOI candidate records from twelve queries and a capped citation-graph expansion. These are candidate literature records, not a complete systematic review. Screening and metadata reconciliation remain necessary before making novelty claims. The engineering comparison must cover domain-specific simulation agents, CAD and component retrieval, structural design, controls, and materials workflows. A catalog survey cannot substitute for direct comparisons with the most relevant systems.

\section{Method}
Let a memory record be $r=(n,i,v,b)$, where $n$ is its namespace, $i$ its identifier, $v$ its version, and $b$ its body. The store rejects changes to an existing namespace--identifier--version tuple. The five namespaces contain procedural skills, engineering tools, parts-source records, dataset-source records, and supporting evidence. The parts and dataset namespaces currently describe sources rather than complete inventories of components or datasets.

\subsection{Taxonomy and specialist scope}
The taxonomy defines domains $D$ and subcategories $C_d$ for each domain $d$. A tool record has explicit catalog memberships $M(r)$. For a requested domain, the scoped tool set is
\begin{equation}
T_d = \{r\in T : M(r)\cap C_d\ne\varnothing\}.
\end{equation}
Selecting one subcategory replaces $C_d$ with that singleton. Unknown domains and categories belonging to another domain are rejected. Membership is editorial metadata; it does not certify task suitability, numerical accuracy, or availability of a tool.

Every domain receives a specialist role specification containing its complete subcategory definitions and instructions to preserve assumptions, units, versions, interfaces, uncertainty, and cited findings. These are profiles for recursive calls, not 71 independently trained or empirically qualified agents. Cross-domain work requires explicit handoffs. The runtime supplies the profiles to the recursive model; automatic optimal routing is not established.

The scope operation selects supporting evidence referenced by the retained tool records and shared parts and dataset records. Missing evidence identifiers are exposed. Skills and shared parts and dataset records remain available. Consequently, domain scoping reduces tool evidence but does not promise that every retained record is relevant to the task.

\subsection{Recursive execution and resource bounds}
The experimental adapter uses the authors' RLM implementation with an external Python dictionary in a Docker environment. The model can inspect this dictionary and request recursive subcalls. The default configuration caps depth at three, iterations at twenty, elapsed execution at 300 seconds, and tokens at 100,000, with four concurrent requests. A dollar threshold is a soft budget guard rather than a guarantee against an in-flight request exceeding the threshold.

The current adapter materializes the selected corpus in host memory before execution. A configurable serialized-character guard limits the exposed payload. It is therefore neither a paged storage service nor an unlimited context window. Each model call remains bounded by the underlying model context, and the system remains bounded by memory, latency, and compute. Scalability must be measured across corpus sizes and task families.

\section{Current implementation evidence}
The local snapshot contains 1,193 skill and reference documents, 4,795 tool records, 48,694 evidence records, 18 parts-source records, and 10 dataset-source records. The tool taxonomy contains 71 domains, 609 subcategories, and 5,658 tool--subcategory membership edges. Counts describe coverage of the imported snapshot, not unique engineering capabilities or validated workflows.

Fourteen deterministic extension tests and ten repository metadata tests passed. The extension tests include immutable addressing, import and retrieval behavior, complete taxonomy counts, scope membership, cross-domain rejection, and missing-evidence reporting. Six local catalog scope audits, covering CAD, optics, CFD, structures, digital electronics, and controls, found no missing referenced evidence IDs in the scoped records. This establishes reference availability for those audits, not factual correctness of the evidence.

A previously recorded synthetic smoke test asked the recursive runtime to combine two load records containing 17 and 29. It returned 46 with both record identifiers. The trace recorded two recursive subcalls, five total model calls, 13,965 input tokens, and 13,517 output tokens over approximately 269 seconds. This test demonstrates one successful recursive execution path. It predates the taxonomy integration and does not validate engineering performance or the specialist profiles.

\section{Planned evaluation}
The primary outcome will be independently scored engineering task correctness. The proposed task families cover CAD, CFD, structural mechanics, controls, optics, and digital electronics. Tasks must specify source artifacts, units, tolerances, compatibility constraints, and independent expected outcomes. Problem-family splits will separate related designs to reduce answer leakage. Answers and scorer artifacts must remain outside model-accessible memory.

The planned conditions are the same model without external resources, progressive skill disclosure, flat retrieval over the memory stores, recursive access to skills, recursive access to all stores, and the full taxonomy-scoped configuration. Models, prompts, and resource ceilings will be matched. Leave-one-store-out and taxonomy-disabled controls will test whether apparent improvements arise from additional evidence or from the scoping mechanism.

The initial planning target is twenty tasks per domain with three repeated runs per condition. A pilot will estimate variance and cost before freezing the final held-out design. Paired outcomes will be summarized with confidence intervals obtained by resampling problem families, with repeated runs aggregated within task. We will report domain-specific correctness, citation support, abstention, constraint violations, latency, token usage, recursion depth, cost, and all failed attempts. No comparative results are available in this draft.

\section{Limitations}
The current corpus is uneven across domains and includes descriptive catalog records that do not establish executable access. Lexical retrieval may miss relevant synonyms. Editorial category errors can exclude useful tools, while shared records can retain irrelevant material. Source provenance does not establish source accuracy. Recursive calls can propagate incorrect assumptions or spend substantial computation on unproductive decomposition. The observed smoke test provides no estimate of robustness, statistical uncertainty, or engineering validity.

The system extends upstream procedural skills and an existing recursive runtime. Establishing novelty requires a completed comparison with adjacent engineering systems and a clear empirical distinction from simply combining retrieval with recursion. Claims about superior performance, unlimited context, autonomous design qualification, or broad domain competence are not supported by the present evidence.

\section*{AI use statement}
Generative AI assisted with the conceptual framework, implementation, literature discovery and synthesis, taxonomy integration, experimental protocol, and manuscript drafting. The synthetic recursion fixture was used for software validation rather than as a scientific dataset. Human review of the manuscript, citations, experimental design, and implementation remains pending. Before submission, the authors must verify the AI-assisted artifacts and update this statement to accurately describe completed review and their responsibility for the final work.

\section*{Ethics statement}
Engineering recommendations can fail when units, constraints, or source assumptions are wrong. Catalog records and generated recommendations require independent validation before consequential use. Redistributed artifacts must respect source licenses and exclude credentials or private catalog contents. This draft reports no human-subject study.

\section*{Reproducibility statement}
The code includes deterministic tests and a recorded synthetic recursive execution trace. The taxonomy appendix is generated from the same versioned definitions used by the runtime. The accompanying evaluation protocol specifies run metadata and artifact requirements. A fully anonymous reproducibility bundle and completed comparative experiment traces remain to be prepared before submission.

\bibliography{references}
\bibliographystyle{iclr2027_conference}
\appendix
\section{Complete engineering taxonomy}
These editorial categories support browsing and scoped retrieval. Membership does not establish tool qualification. Counts refer to the pinned catalog snapshot; overlapping memberships are allowed.
\subsection{Computational Geometry and CAD Kernels}
\noindent Domain: \texttt{geometry}. Primary tool records: 71.
\begin{longtable}{p{0.37\linewidth}p{0.47\linewidth}r}
ID & Subcategory & Count\\\hline\endhead
geometry-solid-modeling-kernels-and-cad-bindings & Solid Modeling Kernels and CAD Bindings & 13\\
geometry-mesh-and-polygon-geometry-processing & Mesh and Polygon Geometry Processing & 23\\
geometry-implicit-volumetric-and-generative-geometry & Implicit Volumetric and Generative Geometry & 6\\
geometry-geometry-translation-and-scene-interchange & Geometry Translation and Scene Interchange & 26\\
geometry-spatial-and-computational-geometry & Spatial and Computational Geometry & 4\\
geometry-point-clouds-and-3d-reconstruction & Point Clouds and 3D Reconstruction & 2\\
geometry-cad-quality-validation-and-healing & CAD Quality Validation and Healing & 3\\
geometry-curve-and-surface-construction & Curve and Surface Construction & 1\\
geometry-geometric-constraints-and-parametric-solvers & Geometric Constraints and Parametric Solvers & 2\\
geometry-engineering-model-visualization-sdks & Engineering Model Visualization SDKs & 1\\
geometry-mesh-geometry-editing-and-reconstruction & Mesh Geometry Editing and Reconstruction & 1\\
\end{longtable}
\subsection{Computer-Aided Design}
\noindent Domain: \texttt{cad}. Primary tool records: 65.
\begin{longtable}{p{0.37\linewidth}p{0.47\linewidth}r}
ID & Subcategory & Count\\\hline\endhead
cad-mechanical & Mechanical and Product Engineering & 25\\
cad-drafting & General Drafting and Technical Documentation & 11\\
cad-architecture & Architectural and Building Design & 7\\
cad-industrial & Industrial Design and Surface Modeling & 12\\
cad-plant & Plant and Factory Layout Design & 3\\
cad-marine & Naval Architecture and Ship Design & 1\\
cad-aeronautical & Aircraft Geometry Design & 1\\
cad-tooling & Mold, Die and Manufacturing Tooling Design & 6\\
cad-computational & Computational and Programmatic Design & 9\\
cad-preparation & Geometry Preparation for Simulation and Manufacturing & 3\\
cad-general-3d & General 3D Geometric Modeling & 4\\
cad-interoperability & CAD File Conversion and Interoperability & 1\\
\end{longtable}
\subsection{Structural Engineering}
\noindent Domain: \texttt{structural}. Primary tool records: 160.
\begin{longtable}{p{0.37\linewidth}p{0.47\linewidth}r}
ID & Subcategory & Count\\\hline\endhead
structural-buildings & Building Structural Design & 17\\
structural-bridges & Bridge and Transportation Structures & 30\\
structural-fea & Structural Analysis and Finite Element Workflows & 90\\
structural-nonlinear & Nonlinear Mechanics, Damage and Fracture & 22\\
structural-dynamics & Structural Dynamics, Seismic and Impact Analysis & 21\\
structural-components & Structural Connections and Component Design & 6\\
structural-composites & Composite and Specialized Structural Materials & 8\\
structural-fatigue & Structural Fatigue and Durability & 4\\
structural-foundations & Structural Foundation Design & 1\\
structural-materials & Structural Material Models and Constitutive Libraries & 2\\
structural-piping & Piping and Equipment Structural Analysis & 2\\
structural-tensile-membranes-cable-nets-and-form-finding & Tensile Membranes Cable Nets and Form Finding & 11\\
structural-construction-stages-and-time-dependent-analysis & Construction Stages and Time-Dependent Analysis & 1\\
structural-prestressed-and-post-tensioned-concrete & Prestressed and Post-Tensioned Concrete & 1\\
structural-composite-modeling-and-results-automation & Composite Modeling and Results Automation & 3\\
structural-composite-and-metallic-structural-sizing & Composite and Metallic Structural Sizing & 1\\
structural-crash-impact-and-explicit-dynamics & Crash Impact and Explicit Dynamics & 1\\
\end{longtable}
\subsection{Building Information Modeling}
\noindent Domain: \texttt{bim}. Primary tool records: 57.
\begin{longtable}{p{0.37\linewidth}p{0.47\linewidth}r}
ID & Subcategory & Count\\\hline\endhead
bim-bim-authoring-and-structural-detailing & BIM Authoring and Structural Detailing & 14\\
bim-bim-coordination-and-model-checking & BIM Coordination and Model Checking & 12\\
bim-bim-interoperability-and-developer-libraries & BIM Interoperability and Developer Libraries & 7\\
bim-bim-collaboration-and-model-viewing & BIM Collaboration and Model Viewing & 18\\
bim-construction-planning-and-4d-sequencing & Construction Planning and 4D Sequencing & 1\\
bim-building-life-cycle-assessment-and-embodied-impacts & Building Life-Cycle Assessment and Embodied Impacts & 3\\
bim-architectural-space-planning-and-design-feasibility & Architectural Space Planning and Design Feasibility & 3\\
bim-mep-design-and-building-services-coordination & MEP Design and Building Services Coordination & 2\\
bim-building-requirements-and-room-data-management & Building Requirements and Room Data Management & 1\\
bim-facility-digital-twins-and-asset-information & Facility Digital Twins and Asset Information & 1\\
bim-civil-infrastructure-conceptual-design & Civil Infrastructure Conceptual Design & 1\\
\end{longtable}
\subsection{Geotechnical Engineering}
\noindent Domain: \texttt{geotechnical}. Primary tool records: 70.
\begin{longtable}{p{0.37\linewidth}p{0.47\linewidth}r}
ID & Subcategory & Count\\\hline\endhead
geotechnical-geotechnical-finite-element-and-continuum-analysis & Geotechnical Finite-Element and Continuum Analysis & 21\\
geotechnical-slope-stability-and-groundwater-seepage & Slope Stability and Groundwater Seepage & 24\\
geotechnical-foundations-retaining-structures-and-excavations & Foundations Retaining Structures and Excavations & 23\\
geotechnical-rock-mechanics-tunneling-and-discontinuities & Rock Mechanics Tunneling and Discontinuities & 15\\
geotechnical-geotechnical-investigation-and-ground-models & Geotechnical Investigation and Ground Models & 1\\
geotechnical-seismic-site-response-and-liquefaction & Seismic Site Response and Liquefaction & 2\\
geotechnical-discrete-element-and-granular-geomechanics & Discrete-Element and Granular Geomechanics & 1\\
geotechnical-mine-geomechanics-and-material-flow & Mine Geomechanics and Material Flow & 2\\
geotechnical-reinforced-soils-and-geosynthetic-structures & Reinforced Soils and Geosynthetic Structures & 2\\
\end{longtable}
\subsection{Computational Fluid Dynamics}
\noindent Domain: \texttt{cfd}. Primary tool records: 131.
\begin{longtable}{p{0.37\linewidth}p{0.47\linewidth}r}
ID & Subcategory & Count\\\hline\endhead
cfd-aerodynamics & Aerodynamics and External Flow & 10\\
cfd-multiphase & Multiphase and Free-Surface Flow & 19\\
cfd-combustion & Combustion and Reacting Flow & 20\\
cfd-compressible & Compressible Flow and Gas Dynamics & 13\\
cfd-particles & Particle and Kinetic Flow Methods & 29\\
cfd-thermal & Thermal Fluids and Conjugate Heat Transfer & 20\\
cfd-turbomachinery & Turbomachinery and Rotating Flow & 5\\
cfd-general & CFD Solvers and Flow Modeling Frameworks & 49\\
cfd-marine & Marine and Coastal Hydrodynamics & 9\\
cfd-fsi & Fluid--Structure Interaction & 1\\
cfd-fire & Fire Simulation and Evacuation Workflows & 5\\
cfd-electrohydrodynamics-and-coupled-field-flows & Electrohydrodynamics and Coupled Field Flows & 1\\
\end{longtable}
\subsection{Radio Frequency and Electromagnetic Engineering}
\noindent Domain: \texttt{rf-em}. Primary tool records: 96.
\begin{longtable}{p{0.37\linewidth}p{0.47\linewidth}r}
ID & Subcategory & Count\\\hline\endhead
rf-em-antennas & Antenna Design and Electromagnetic Scattering & 14\\
rf-em-microwave & Microwave, RF and High-Frequency Circuit Design & 20\\
rf-em-compatibility & Electromagnetic Compatibility and Interference & 6\\
rf-em-low-frequency & Low-Frequency Fields and Electromechanics & 16\\
rf-em-fullwave & Full-Wave Electromagnetic Simulation & 22\\
rf-em-interconnect & Signal Integrity and Electronic Interconnects & 16\\
rf-em-quantum & Quantum Transport and Electronic Device Modeling & 15\\
rf-em-superconducting & Superconducting and Josephson Circuit Modeling & 7\\
rf-em-layout & Microelectronic and Photonic Layout Libraries & 3\\
rf-em-photonics & Optical and Photonic Electromagnetic Modeling & 2\\
rf-em-propagation & Wireless Propagation and Network Planning & 1\\
\end{longtable}
\subsection{Optical Engineering and Photonics}
\noindent Domain: \texttt{optics}. Primary tool records: 167.
\begin{longtable}{p{0.37\linewidth}p{0.47\linewidth}r}
ID & Subcategory & Count\\\hline\endhead
optics-imaging & Optical System Design and Imaging & 33\\
optics-illumination & Illumination and Nonimaging Optics & 17\\
optics-wave & Wave Optics and Electromagnetic Propagation & 52\\
optics-integrated & Integrated Photonics and Photonic Circuits & 13\\
optics-laser & Lasers, Ultrafast and Nonlinear Optics & 35\\
optics-fiber & Fiber Optics and Optical Communications & 17\\
optics-coatings & Optical Coatings and Thin Films & 2\\
optics-characterization & Optical Measurement and Characterization & 5\\
optics-devices & Optoelectronic and Semiconductor Device Modeling & 1\\
optics-inverse & Optical Optimization and Inverse Design & 10\\
optics-adaptive & Adaptive Optics and Wavefront Control & 10\\
optics-optomechanical-and-thermal-performance-integration & Optomechanical and Thermal Performance Integration & 5\\
\end{longtable}
\subsection{Integrated Circuit Design and Semiconductor Devices}
\noindent Domain: \texttt{eda-digital}. Primary tool records: 55.
\begin{longtable}{p{0.37\linewidth}p{0.47\linewidth}r}
ID & Subcategory & Count\\\hline\endhead
eda-digital-logic-synthesis-and-physical-implementation & Logic Synthesis and Physical Implementation & 13\\
eda-digital-timing-power-and-signoff-analysis & Timing Power and Signoff Analysis & 8\\
eda-digital-digital-simulation-and-verification & Digital Simulation and Verification & 27\\
eda-digital-test-design-and-silicon-debug & Test Design and Silicon Debug & 6\\
eda-digital-custom-layout-and-physical-verification & Custom Layout and Physical Verification & 3\\
eda-digital-eda-flow-orchestration-and-build-infrastructure & EDA Flow Orchestration and Build Infrastructure & 2\\
eda-digital-microfabrication-and-mask-layout & Microfabrication and Mask Layout & 1\\
eda-digital-semiconductor-device-simulation & Semiconductor Device Simulation & 1\\
\end{longtable}
\subsection{FPGA Design and Verification}
\noindent Domain: \texttt{fpga}. Primary tool records: 34.
\begin{longtable}{p{0.37\linewidth}p{0.47\linewidth}r}
ID & Subcategory & Count\\\hline\endhead
fpga-fpga-synthesis-and-place-and-route & FPGA Synthesis and Place-and-Route & 19\\
fpga-high-level-synthesis-and-embedded-software & High-Level Synthesis and Embedded Software & 4\\
fpga-fpga-simulation-verification-and-debug & FPGA Simulation Verification and Debug & 12\\
fpga-fpga-programming-and-hardware-integration & FPGA Programming and Hardware Integration & 12\\
\end{longtable}
\subsection{Electronic Circuit Simulation}
\noindent Domain: \texttt{spice}. Primary tool records: 39.
\begin{longtable}{p{0.37\linewidth}p{0.47\linewidth}r}
ID & Subcategory & Count\\\hline\endhead
spice-spice-and-analog-circuit-simulation & SPICE and Analog Circuit Simulation & 31\\
spice-power-electronics-and-switching-systems & Power Electronics and Switching Systems & 3\\
spice-mixed-signal-and-behavioral-modeling & Mixed-Signal and Behavioral Modeling & 4\\
spice-circuit-design-automation-and-schematic-interfaces & Circuit Design Automation and Schematic Interfaces & 7\\
spice-rf-and-microwave-circuit-simulation & RF and Microwave Circuit Simulation & 4\\
\end{longtable}
\subsection{Manufacturing Engineering}
\noindent Domain: \texttt{manufacturing}. Primary tool records: 203.
\begin{longtable}{p{0.37\linewidth}p{0.47\linewidth}r}
ID & Subcategory & Count\\\hline\endhead
manufacturing-machining & CNC Machining and CAM & 62\\
manufacturing-additive & Additive Manufacturing & 23\\
manufacturing-forming & Metal Forming and Sheet Metal Processing & 22\\
manufacturing-molding & Casting and Polymer Processing & 22\\
manufacturing-joining & Welding and Joining Processes & 6\\
manufacturing-thermal & Heat Treatment and Surface Processing & 14\\
manufacturing-planning & Manufacturing Process Planning and Operations & 4\\
manufacturing-composites & Composite Manufacturing & 16\\
manufacturing-simulation & Manufacturing Process Simulation and Optimization & 13\\
manufacturing-preparation & Manufacturing Geometry and Job Preparation & 2\\
manufacturing-response & Materials Response and Microstructure for Manufacturing & 3\\
manufacturing-fluid & Process Fluid Simulation & 1\\
manufacturing-polymer-processing-and-extrusion-simulation & Polymer Processing and Extrusion Simulation & 1\\
manufacturing-electronics-encapsulation-and-packaging-processes & Electronics Encapsulation and Packaging Processes & 2\\
manufacturing-composite-forming-draping-and-layup-preparation & Composite Forming Draping and Layup Preparation & 4\\
manufacturing-metal-forming-and-process-simulation & Metal Forming and Process Simulation & 8\\
manufacturing-heat-treatment-and-microstructure-evolution & Heat Treatment and Microstructure Evolution & 2\\
manufacturing-filament-winding-and-composite-toolpaths & Filament Winding and Composite Toolpaths & 2\\
manufacturing-manufacturing-process-planning & Manufacturing Process Planning & 1\\
manufacturing-textile-architecture-weaving-and-fabric-simulation & Textile Architecture Weaving and Fabric Simulation & 10\\
manufacturing-membrane-cutting-patterns-and-fabrication & Membrane Cutting Patterns and Fabrication & 2\\
manufacturing-composite-infusion-and-permeability & Composite Infusion and Permeability & 2\\
manufacturing-casting-and-solidification & Casting and Solidification & 3\\
manufacturing-automated-fiber-placement-and-tape-laying & Automated Fiber Placement and Tape Laying & 6\\
manufacturing-composite-cure-and-process-induced-distortion & Composite Cure and Process-Induced Distortion & 1\\
manufacturing-composite-layup-design-and-manufacturing-preparation & Composite Layup Design and Manufacturing Preparation & 2\\
manufacturing-composite-material-tracking-and-production-management & Composite Material Tracking and Production Management & 1\\
manufacturing-sheet-cutting-nesting-and-cnc-programming & Sheet Cutting Nesting and CNC Programming & 1\\
\end{longtable}
\subsection{Materials Science and Engineering}
\noindent Domain: \texttt{materials}. Primary tool records: 96.
\begin{longtable}{p{0.37\linewidth}p{0.47\linewidth}r}
ID & Subcategory & Count\\\hline\endhead
materials-thermodynamics & Materials Thermodynamics and Phase Equilibria & 7\\
materials-microstructure & Microstructure and Phase-Field Modeling & 27\\
materials-constitutive & Constitutive Modeling and Material Response & 10\\
materials-processing & Materials Processing and Heat Treatment & 5\\
materials-crystallography & Crystallography, Diffraction and Scattering & 27\\
materials-magnetism & Micromagnetics and Magnetic Materials & 13\\
materials-discovery & Materials Discovery and Informatics & 3\\
materials-composites & Composite and Multiscale Materials & 2\\
materials-kinetics & Surface Chemistry and Kinetic Modeling & 4\\
materials-data & Materials Data, Structure Preparation and Characterization & 4\\
materials-potentials & Interatomic Potentials and Atomistic Material Models & 1\\
materials-composite-micromechanics-and-homogenization & Composite Micromechanics and Homogenization & 2\\
materials-thermal-protection-and-ablative-material-response & Thermal Protection and Ablative Material Response & 1\\
\end{longtable}
\subsection{Robotics Engineering}
\noindent Domain: \texttt{robotics}. Primary tool records: 68.
\begin{longtable}{p{0.37\linewidth}p{0.47\linewidth}r}
ID & Subcategory & Count\\\hline\endhead
robotics-programming & Industrial Robot Programming and Workcells & 27\\
robotics-simulation & Robot Simulation and Virtual Environments & 21\\
robotics-dynamics & Robot Kinematics, Dynamics and Contact Physics & 16\\
robotics-planning & Motion Planning and Trajectory Optimization & 13\\
robotics-learning & Robot Learning and Synthetic Data & 2\\
robotics-middleware & Robotics Middleware and Application Frameworks & 4\\
robotics-commissioning & Manufacturing Simulation and Virtual Commissioning & 7\\
robotics-education & Robotics Education and Algorithm Prototyping & 1\\
robotics-small-farm & Small-Farm Automation & 1\\
robotics-fleet-coordination-and-facility-integration & Fleet Coordination and Facility Integration & 1\\
\end{longtable}
\subsection{Electrical Power Systems Engineering}
\noindent Domain: \texttt{power}. Primary tool records: 51.
\begin{longtable}{p{0.37\linewidth}p{0.47\linewidth}r}
ID & Subcategory & Count\\\hline\endhead
power-power-flow-distribution-and-network-planning & Power Flow Distribution and Network Planning & 24\\
power-electromagnetic-transients-and-real-time-simulation & Electromagnetic Transients and Real-Time Simulation & 9\\
power-protection-fault-and-electrical-safety-analysis & Protection Fault and Electrical Safety Analysis & 12\\
power-grid-dynamics-stability-and-frequency-control & Grid Dynamics Stability and Frequency Control & 7\\
power-electric-machine-and-motor-design & Electric Machine and Motor Design & 1\\
power-energy-markets-resource-adequacy-and-capacity-expansion & Energy Markets Resource Adequacy and Capacity Expansion & 13\\
power-railway-traction-power-and-electrification & Railway Traction Power and Electrification & 2\\
\end{longtable}
\subsection{Scientific Machine Learning}
\noindent Domain: \texttt{sciml}. Primary tool records: 49.
\begin{longtable}{p{0.37\linewidth}p{0.47\linewidth}r}
ID & Subcategory & Count\\\hline\endhead
sciml-physics-informed-neural-networks & Physics-Informed Neural Networks & 6\\
sciml-neural-operators-and-learned-surrogates & Neural Operators and Learned Surrogates & 10\\
sciml-differentiable-simulation-and-scientific-solvers & Differentiable Simulation and Scientific Solvers & 26\\
sciml-machine-learning-frameworks-and-regression & Machine Learning Frameworks and Regression & 13\\
sciml-symbolic-regression-and-scientific-programming & Symbolic Regression and Scientific Programming & 2\\
\end{longtable}
\subsection{Laboratory Automation and Instrumentation}
\noindent Domain: \texttt{lab}. Primary tool records: 69.
\begin{longtable}{p{0.37\linewidth}p{0.47\linewidth}r}
ID & Subcategory & Count\\\hline\endhead
lab-control & Instrument Control and Data Acquisition & 52\\
lab-automation & Experiment Orchestration and Laboratory Automation & 14\\
lab-processing & Experimental Data Processing and Analysis & 3\\
lab-integration & Instrument Communication and Integration Libraries & 12\\
\end{longtable}
\subsection{Building Energy and HVAC Engineering}
\noindent Domain: \texttt{building-energy}. Primary tool records: 56.
\begin{longtable}{p{0.37\linewidth}p{0.47\linewidth}r}
ID & Subcategory & Count\\\hline\endhead
building-energy-whole-building-energy-simulation & Whole-Building Energy Simulation & 15\\
building-energy-district-and-urban-energy-modeling & District and Urban Energy Modeling & 9\\
building-energy-hvac-systems-and-building-controls & HVAC Systems and Building Controls & 27\\
building-energy-envelope-moisture-and-passive-design & Envelope Moisture and Passive Design & 8\\
building-energy-daylighting-and-early-design-assessment & Daylighting and Early Design Assessment & 8\\
building-energy-energy-code-compliance-and-ratings & Energy Code Compliance and Ratings & 3\\
building-energy-building-airflow-and-indoor-air-quality & Building Airflow and Indoor Air Quality & 1\\
\end{longtable}
\subsection{Water Resources and Environmental Modeling}
\noindent Domain: \texttt{water}. Primary tool records: 112.
\begin{longtable}{p{0.37\linewidth}p{0.47\linewidth}r}
ID & Subcategory & Count\\\hline\endhead
water-hydrology & Watershed Hydrology and Rainfall--Runoff Modeling & 29\\
water-rivers & River Hydraulics and Flood Modeling & 34\\
water-drainage & Urban Drainage and Sewer Systems & 7\\
water-distribution & Water Distribution and Pipe Networks & 10\\
water-quality & Water Quality and Aquatic Ecosystems & 11\\
water-coastal & Coastal, Estuarine and Ocean Hydrodynamics & 34\\
water-groundwater & Groundwater and Integrated Water Resources & 19\\
water-sediment & Sediment Transport and Morphodynamics & 5\\
water-hydraulic-model-preparation-and-visualization & Hydraulic Model Preparation and Visualization & 1\\
water-hydrological-model-preparation-and-gis-integration & Hydrological Model Preparation and GIS Integration & 1\\
water-hydrologic-frequency-and-dam-risk-analysis & Hydrologic Frequency and Dam-Risk Analysis & 2\\
\end{longtable}
\subsection{Electrochemical Engineering}
\noindent Domain: \texttt{electrochemistry}. Primary tool records: 59.
\begin{longtable}{p{0.37\linewidth}p{0.47\linewidth}r}
ID & Subcategory & Count\\\hline\endhead
electrochemistry-battery-cell-and-pack-modeling & Battery Cell and Pack Modeling & 29\\
electrochemistry-fuel-cells-electrolyzers-and-hydrogen-systems & Fuel Cells Electrolyzers and Hydrogen Systems & 15\\
electrochemistry-electrochemical-kinetics-and-transport & Electrochemical Kinetics and Transport & 6\\
electrochemistry-electrochemical-instrument-control-and-experiment-analysis & Electrochemical Instrument Control and Experiment Analysis & 1\\
electrochemistry-impedance-spectroscopy-and-equivalent-circuit-analysis & Impedance Spectroscopy and Equivalent-Circuit Analysis & 3\\
electrochemistry-battery-manufacturing-and-cost-analysis & Battery Manufacturing and Cost Analysis & 1\\
electrochemistry-hydrogen-production-delivery-and-refueling-economics & Hydrogen Production Delivery and Refueling Economics & 6\\
electrochemistry-battery-aging-and-lifetime-prediction & Battery Aging and Lifetime Prediction & 1\\
\end{longtable}
\subsection{Chemical and Process Engineering}
\noindent Domain: \texttt{chemical}. Primary tool records: 69.
\begin{longtable}{p{0.37\linewidth}p{0.47\linewidth}r}
ID & Subcategory & Count\\\hline\endhead
chemical-process-flowsheeting-and-plant-simulation & Process Flowsheeting and Plant Simulation & 34\\
chemical-reaction-kinetics-and-reactor-modeling & Reaction Kinetics and Reactor Modeling & 15\\
chemical-thermophysical-properties-and-transport & Thermophysical Properties and Transport & 16\\
chemical-separation-processes-and-unit-operations & Separation Processes and Unit Operations & 6\\
chemical-batch-and-particulate-process-engineering & Batch and Particulate Process Engineering & 7\\
chemical-chemical-structure-representation-and-cheminformatics & Chemical Structure Representation and Cheminformatics & 5\\
chemical-mixing-and-agitated-process-equipment & Mixing and Agitated Process Equipment & 1\\
chemical-chemical-kinetics-and-mechanism-reduction & Chemical Kinetics and Mechanism Reduction & 1\\
chemical-chemical-equilibrium-and-thermochemistry & Chemical Equilibrium and Thermochemistry & 1\\
\end{longtable}
\subsection{Porous Media and Subsurface Flow}
\noindent Domain: \texttt{porous-media}. Primary tool records: 76.
\begin{longtable}{p{0.37\linewidth}p{0.47\linewidth}r}
ID & Subcategory & Count\\\hline\endhead
porous-media-reservoir-and-multiphase-flow-simulation & Reservoir and Multiphase Flow Simulation & 22\\
porous-media-reactive-transport-and-geochemistry & Reactive Transport and Geochemistry & 12\\
porous-media-groundwater-solute-and-unsaturated-transport & Groundwater Solute and Unsaturated Transport & 9\\
porous-media-pore-scale-and-porous-media-modeling & Pore-Scale and Porous-Media Modeling & 1\\
porous-media-fractured-media-flow-and-geomechanics & Fractured-Media Flow and Geomechanics & 1\\
porous-media-continuum-porous-media-flow-and-transport & Continuum Porous-Media Flow and Transport & 3\\
porous-media-pore-scale-flow-and-multiphase-transport & Pore-Scale Flow and Multiphase Transport & 2\\
porous-media-reservoir-history-matching-uncertainty-and-optimization & Reservoir History Matching Uncertainty and Optimization & 1\\
porous-media-reservoir-wellbore-and-production-system-coupling & Reservoir Wellbore and Production-System Coupling & 5\\
porous-media-subsurface-engineering-platforms-and-data-integration & Subsurface Engineering Platforms and Data Integration & 1\\
porous-media-drilling-well-control-and-wellbore-hydraulics & Drilling Well Control and Wellbore Hydraulics & 1\\
porous-media-geological-carbon-storage-assessment & Geological Carbon Storage Assessment & 1\\
porous-media-hydraulic-fracturing-and-well-stimulation & Hydraulic Fracturing and Well Stimulation & 5\\
porous-media-well-testing-production-diagnostics-and-surveillance & Well Testing Production Diagnostics and Surveillance & 5\\
porous-media-well-design-completion-and-production-hydraulics & Well Design Completion and Production Hydraulics & 3\\
porous-media-well-design-and-integrity & Well Design and Integrity & 4\\
porous-media-reservoir-model-automation-and-interfaces & Reservoir Model Automation and Interfaces & 1\\
porous-media-well-lifecycle-data-and-operations & Well Lifecycle Data and Operations & 2\\
porous-media-reservoir-model-preparation-and-visualization & Reservoir Model Preparation and Visualization & 2\\
\end{longtable}
\subsection{Multiphysics Modeling and Simulation}
\noindent Domain: \texttt{multiphysics}. Primary tool records: 51.
\begin{longtable}{p{0.37\linewidth}p{0.47\linewidth}r}
ID & Subcategory & Count\\\hline\endhead
multiphysics-pde-and-finite-element-frameworks & PDE and Finite-Element Frameworks & 28\\
multiphysics-multiphysics-simulation-suites & Multiphysics Simulation Suites & 8\\
multiphysics-coupling-and-co-simulation-infrastructure & Coupling and Co-Simulation Infrastructure & 3\\
multiphysics-multibody-and-mechanism-simulation & Multibody and Mechanism Simulation & 6\\
multiphysics-thermal-and-electromagnetic-systems & Thermal and Electromagnetic Systems & 5\\
multiphysics-flexible-rods-cables-and-slender-body-mechanics & Flexible Rods Cables and Slender-Body Mechanics & 1\\
multiphysics-discrete-element-and-granular-mechanics & Discrete-Element and Granular Mechanics & 2\\
\end{longtable}
\subsection{Acoustics and Vibration Engineering}
\noindent Domain: \texttt{acoustics}. Primary tool records: 57.
\begin{longtable}{p{0.37\linewidth}p{0.47\linewidth}r}
ID & Subcategory & Count\\\hline\endhead
acoustics-room-and-architectural-acoustics & Room and Architectural Acoustics & 17\\
acoustics-environmental-noise-and-sound-propagation & Environmental Noise and Sound Propagation & 4\\
acoustics-vibroacoustics-and-structural-sound & Vibroacoustics and Structural Sound & 6\\
acoustics-acoustic-wave-solvers-and-signal-analysis & Acoustic Wave Solvers and Signal Analysis & 10\\
acoustics-absorbers-and-acoustic-materials & Absorbers and Acoustic Materials & 3\\
acoustics-electroacoustic-systems-and-loudspeakers & Electroacoustic Systems and Loudspeakers & 2\\
acoustics-underwater-acoustics-and-sonar & Underwater Acoustics and Sonar & 10\\
acoustics-acoustic-measurement-and-source-localization & Acoustic Measurement and Source Localization & 6\\
acoustics-thermoacoustics-and-combustion-instability & Thermoacoustics and Combustion Instability & 1\\
acoustics-ultrasound-and-piezoelectric-transducers & Ultrasound and Piezoelectric Transducers & 1\\
\end{longtable}
\subsection{Renewable Energy Engineering}
\noindent Domain: \texttt{renewables}. Primary tool records: 50.
\begin{longtable}{p{0.37\linewidth}p{0.47\linewidth}r}
ID & Subcategory & Count\\\hline\endhead
renewables-wind-turbine-aerodynamics-and-structural-dynamics & Wind Turbine Aerodynamics and Structural Dynamics & 6\\
renewables-wind-resource-assessment-and-farm-layout & Wind Resource Assessment and Farm Layout & 10\\
renewables-solar-photovoltaic-design-and-yield & Solar Photovoltaic Design and Yield & 14\\
renewables-marine-renewable-energy-systems & Marine Renewable Energy Systems & 8\\
renewables-energy-system-planning-and-technoeconomic-analysis & Energy System Planning and Technoeconomic Analysis & 10\\
renewables-energy-storage-and-distributed-energy-systems & Energy Storage and Distributed Energy Systems & 7\\
\end{longtable}
\subsection{Aeronautical Engineering}
\noindent Domain: \texttt{aerospace}. Primary tool records: 122.
\begin{longtable}{p{0.37\linewidth}p{0.47\linewidth}r}
ID & Subcategory & Count\\\hline\endhead
aerospace-aircraft-conceptual-design-and-sizing & Aircraft Conceptual Design and Sizing & 20\\
aerospace-aerodynamics-and-aeroelasticity & Aerodynamics and Aeroelasticity & 25\\
aerospace-flight-dynamics-and-simulation & Flight Dynamics and Simulation & 12\\
aerospace-propulsion-and-engine-performance & Propulsion and Engine Performance & 50\\
aerospace-rotorcraft-and-rotor-aerodynamics & Rotorcraft and Rotor Aerodynamics & 9\\
aerospace-aeroacoustics-and-propulsion-noise & Aeroacoustics and Propulsion Noise & 10\\
aerospace-flight-control-and-handling-qualities & Flight Control and Handling Qualities & 13\\
aerospace-rocket-design-and-flight-simulation & Rocket Design and Flight Simulation & 4\\
aerospace-aircraft-performance-and-mission-analysis & Aircraft Performance and Mission Analysis & 3\\
aerospace-open-aircraft-designs & Open Aircraft Designs and Build Resources & 1\\
\end{longtable}
\subsection{Space Systems and Astronautical Engineering}
\noindent Domain: \texttt{space}. Primary tool records: 97.
\begin{longtable}{p{0.37\linewidth}p{0.47\linewidth}r}
ID & Subcategory & Count\\\hline\endhead
space-orbital-mechanics-and-trajectory-design & Orbital Mechanics and Trajectory Design & 42\\
space-spacecraft-dynamics-guidance-and-control & Spacecraft Dynamics Guidance and Control & 9\\
space-spacecraft-thermal-and-environmental-analysis & Spacecraft Thermal and Environmental Analysis & 28\\
space-mission-systems-communications-and-coverage & Mission Systems Communications and Coverage & 9\\
space-mission-operations-and-flight-software-testing & Mission Operations and Flight Software Testing & 8\\
space-plume-interaction-and-surface-operations & Plume Interaction and Surface Operations & 1\\
space-planetary-surface-operations-and-resource-utilization & Planetary Surface Operations and Resource Utilization & 2\\
space-spacecraft-power-and-energy-systems & Spacecraft Power and Energy Systems & 3\\
space-exoplanet-mission-design-and-yield-analysis & Exoplanet Mission Design and Yield Analysis & 1\\
\end{longtable}
\subsection{Nuclear Science and Engineering}
\noindent Domain: \texttt{nuclear}. Primary tool records: 117.
\begin{longtable}{p{0.37\linewidth}p{0.47\linewidth}r}
ID & Subcategory & Count\\\hline\endhead
nuclear-transport & Radiation Transport and Shielding & 29\\
nuclear-reactor & Reactor Physics and Core Analysis & 44\\
nuclear-thermal & Nuclear Thermal Hydraulics and Safety Analysis & 19\\
nuclear-fuel & Nuclear Fuel Performance and Fuel Cycles & 30\\
nuclear-fusion & Fusion Engineering and Neutronics & 2\\
nuclear-environment & Environmental Radiological Assessment and Dosimetry & 7\\
nuclear-risk & Nuclear Reliability and Probabilistic Risk Assessment & 1\\
nuclear-integration & Nuclear Analysis Integration and Workflow & 8\\
nuclear-nuclear-data-processing-and-evaluation & Nuclear Data Processing and Evaluation & 6\\
nuclear-transport-benchmarking-and-nuclear-data-validation & Transport Benchmarking and Nuclear-Data Validation & 1\\
nuclear-neutron-instrument-and-scattering-simulation & Neutron Instrument and Scattering Simulation & 2\\
nuclear-radiation-transport-geometry-and-model-preparation & Radiation Transport Geometry and Model Preparation & 1\\
\end{longtable}
\subsection{Numerical Methods and Scientific Computing}
\noindent Domain: \texttt{numerical}. Primary tool records: 90.
\begin{longtable}{p{0.37\linewidth}p{0.47\linewidth}r}
ID & Subcategory & Count\\\hline\endhead
numerical-linear & Linear Algebra and Matrix Computation & 22\\
numerical-differential & Differential Equations and Numerical Solvers & 8\\
numerical-symbolic & Symbolic and Computer Algebra & 11\\
numerical-arrays & Numerical Arrays and Scientific Computing & 14\\
numerical-transforms & Transforms, Approximation and Numerical Analysis & 4\\
numerical-acceleration & Parallel and Accelerated Numerical Computing & 16\\
numerical-precision & Arbitrary-Precision and Interval Arithmetic & 7\\
numerical-optimization & Numerical Optimization and Mathematical Programming & 7\\
numerical-differentiation & Automatic Differentiation & 7\\
numerical-compilation & Scientific Compilation and Code Acceleration & 2\\
numerical-uncertainty & Uncertainty Quantification and Propagation & 2\\
numerical-storage & Scientific Data Storage and Interchange & 1\\
numerical-units & Physical Units and Dimensional Computation & 1\\
numerical-formal-mathematics-theorem-proving-and-smt & Formal Mathematics Theorem Proving and SMT & 7\\
\end{longtable}
\subsection{Scientific Data Analysis and Statistics}
\noindent Domain: \texttt{data-analysis}. Primary tool records: 71.
\begin{longtable}{p{0.37\linewidth}p{0.47\linewidth}r}
ID & Subcategory & Count\\\hline\endhead
data-analysis-tabular-and-multidimensional-data-processing & Tabular and Multidimensional Data Processing & 23\\
data-analysis-scientific-storage-and-data-interchange & Scientific Storage and Data Interchange & 29\\
data-analysis-statistics-and-bayesian-inference & Statistics and Bayesian Inference & 22\\
data-analysis-machine-learning-and-model-interpretation & Machine Learning and Model Interpretation & 3\\
data-analysis-distributed-and-large-scale-data-computing & Distributed and Large-Scale Data Computing & 2\\
data-analysis-time-series-analysis-and-forecasting & Time-Series Analysis and Forecasting & 4\\
data-analysis-time-series-and-forecasting & Time Series and Forecasting & 1\\
\end{longtable}
\subsection{Mathematical Optimization}
\noindent Domain: \texttt{optimization}. Primary tool records: 91.
\begin{longtable}{p{0.37\linewidth}p{0.47\linewidth}r}
ID & Subcategory & Count\\\hline\endhead
optimization-mathematical-programming-and-constraint-solvers & Mathematical Programming and Constraint Solvers & 27\\
optimization-engineering-design-and-multidisciplinary-optimization & Engineering Design and Multidisciplinary Optimization & 24\\
optimization-global-evolutionary-and-black-box-optimization & Global Evolutionary and Black-Box Optimization & 16\\
optimization-uncertainty-sensitivity-and-experimental-design & Uncertainty Sensitivity and Experimental Design & 16\\
optimization-surrogate-modeling-and-parameter-fitting & Surrogate Modeling and Parameter Fitting & 12\\
optimization-optimal-control-and-differentiable-optimization & Optimal Control and Differentiable Optimization & 7\\
optimization-design-study-and-optimization-workflow-orchestration & Design Study and Optimization Workflow Orchestration & 2\\
\end{longtable}
\subsection{Mesh Generation and Processing}
\noindent Domain: \texttt{meshing}. Primary tool records: 48.
\begin{longtable}{p{0.37\linewidth}p{0.47\linewidth}r}
ID & Subcategory & Count\\\hline\endhead
meshing-finite-element-and-volume-mesh-generation & Finite-Element and Volume Mesh Generation & 37\\
meshing-mesh-adaptation-optimization-and-repair & Mesh Adaptation Optimization and Repair & 11\\
meshing-mesh-interchange-and-preprocessing & Mesh Interchange and Preprocessing & 2\\
meshing-specialized-cfd-and-boundary-layer-meshing & Specialized CFD and Boundary-Layer Meshing & 6\\
meshing-thermal-model-geometry-preparation & Thermal Model Geometry Preparation & 1\\
\end{longtable}
\subsection{Printed Circuit Board Design}
\noindent Domain: \texttt{eda-pcb}. Primary tool records: 66.
\begin{longtable}{p{0.37\linewidth}p{0.47\linewidth}r}
ID & Subcategory & Count\\\hline\endhead
eda-pcb-pcb-schematic-and-layout-design & PCB Schematic and Layout Design & 26\\
eda-pcb-signal-power-and-electromagnetic-integrity & Signal Power and Electromagnetic Integrity & 13\\
eda-pcb-advanced-packaging-and-heterogeneous-integration & Advanced Packaging and Heterogeneous Integration & 8\\
eda-pcb-pcb-manufacturing-and-assembly-verification & PCB Manufacturing and Assembly Verification & 17\\
eda-pcb-electronics-reliability-and-thermal-analysis & Electronics Reliability and Thermal Analysis & 2\\
eda-pcb-pcb-stackup-design-and-engineering-calculators & PCB Stackup Design and Engineering Calculators & 5\\
eda-pcb-electronics-thermal-management & Electronics Thermal Management & 5\\
eda-pcb-behavioral-i-o-and-ibis-model-development & Behavioral I/O and IBIS Model Development & 1\\
eda-pcb-package-thermomechanical-reliability & Package Thermomechanical Reliability & 2\\
eda-pcb-electronic-design-governance-and-collaboration & Electronic Design Governance and Collaboration & 1\\
\end{longtable}
\subsection{Control Systems and Dynamic Simulation}
\noindent Domain: \texttt{controls}. Primary tool records: 140.
\begin{longtable}{p{0.37\linewidth}p{0.47\linewidth}r}
ID & Subcategory & Count\\\hline\endhead
controls-physical-system-and-modelica-simulation & Physical System and Modelica Simulation & 25\\
controls-control-design-and-optimal-control & Control Design and Optimal Control & 17\\
controls-co-simulation-and-model-exchange & Co-Simulation and Model Exchange & 11\\
controls-real-time-simulation-and-hardware-in-the-loop & Real-Time Simulation and Hardware-in-the-Loop & 20\\
controls-system-identification-and-state-estimation & System Identification and State Estimation & 6\\
controls-state-machines-and-discrete-event-control & State Machines and Discrete-Event Control & 2\\
controls-differential-equations-and-dynamic-system-solvers & Differential Equations and Dynamic System Solvers & 4\\
controls-mechanical-circuit-and-educational-system-simulation & Mechanical Circuit and Educational System Simulation & 1\\
controls-embedded-control-software-and-code-generation & Embedded Control Software and Code Generation & 4\\
controls-vehicle-and-sensor-simulation & Vehicle and Sensor Simulation & 37\\
controls-rail-vehicle-and-train-dynamics & Rail Vehicle and Train Dynamics & 1\\
controls-multibody-dynamics-and-mechanism-simulation & Multibody Dynamics and Mechanism Simulation & 23\\
controls-industrial-monitoring-and-supervisory-control & Industrial Monitoring and Supervisory Control & 3\\
controls-plc-programming-and-motion-control & PLC Programming and Motion Control & 4\\
\end{longtable}
\subsection{Statistical Mechanics}
\noindent Domain: \texttt{statistical-mechanics}. Primary tool records: 44.
\begin{longtable}{p{0.37\linewidth}p{0.47\linewidth}r}
ID & Subcategory & Count\\\hline\endhead
statistical-mechanics-molecular-and-particle-dynamics & Molecular and Particle Dynamics & 10\\
statistical-mechanics-monte-carlo-and-ensemble-sampling & Monte Carlo and Ensemble Sampling & 17\\
statistical-mechanics-enhanced-sampling-and-free-energy & Enhanced Sampling and Free Energy & 11\\
statistical-mechanics-trajectory-analysis-and-statistical-estimation & Trajectory Analysis and Statistical Estimation & 2\\
statistical-mechanics-mesoscopic-and-reaction-diffusion-simulation & Mesoscopic and Reaction-Diffusion Simulation & 2\\
statistical-mechanics-thermodynamics-and-molecular-property-prediction & Thermodynamics and Molecular Property Prediction & 1\\
statistical-mechanics-magnetic-structures-and-spin-dynamics & Magnetic Structures and Spin Dynamics & 4\\
\end{longtable}
\subsection{Atomistic Modeling and Quantum Chemistry}
\noindent Domain: \texttt{atomistic}. Primary tool records: 85.
\begin{longtable}{p{0.37\linewidth}p{0.47\linewidth}r}
ID & Subcategory & Count\\\hline\endhead
atomistic-electronic & Electronic Structure and Density Functional Theory & 42\\
atomistic-quantum-chemistry & Molecular Quantum Chemistry & 20\\
atomistic-molecular-dynamics & Molecular Dynamics and Atomistic Simulation & 15\\
atomistic-potentials & Interatomic Potentials and Machine Learning & 9\\
atomistic-workflows & Atomistic Workflows and Structure Analysis & 21\\
atomistic-transport & Electronic Transport and Mesoscopic Modeling & 4\\
atomistic-excited & Excited-State and Nonadiabatic Dynamics & 1\\
atomistic-vibrations & Lattice Dynamics and Phonons & 1\\
atomistic-integrated & Integrated Molecular and Materials Modeling & 2\\
atomistic-molecular-force-fields-and-parameterization & Molecular Force Fields and Parameterization & 1\\
\end{longtable}
\subsection{Thermodynamics and Thermophysical Properties}
\noindent Domain: \texttt{thermophysical}. Primary tool records: 99.
\begin{longtable}{p{0.37\linewidth}p{0.47\linewidth}r}
ID & Subcategory & Count\\\hline\endhead
thermophysical-fluid-properties-and-equations-of-state & Fluid Properties and Equations of State & 18\\
thermophysical-thermal-hydraulic-networks-and-transients & Thermal Hydraulic Networks and Transients & 39\\
thermophysical-heat-exchangers-and-thermal-equipment & Heat Exchangers and Thermal Equipment & 7\\
thermophysical-thermodynamic-cycles-and-energy-systems & Thermodynamic Cycles and Energy Systems & 10\\
thermophysical-cryogenic-and-vacuum-systems & Cryogenic and Vacuum Systems & 20\\
thermophysical-psychrometrics-and-moist-air-properties & Psychrometrics and Moist-Air Properties & 1\\
thermophysical-heat-transfer-simulation & Heat Transfer Simulation & 4\\
thermophysical-thermal-model-automation-and-interfaces & Thermal Model Automation and Interfaces & 1\\
thermophysical-ground-source-heat-and-geothermal-borefields & Ground-Source Heat and Geothermal Borefields & 2\\
thermophysical-cryogenic-systems-and-superconducting-magnet-thermal-analysis & Cryogenic Systems and Superconducting Magnet Thermal Analysis & 4\\
thermophysical-heat-exchanger-design-and-rating & Heat Exchanger Design and Rating & 1\\
\end{longtable}
\subsection{Scientific Visualization}
\noindent Domain: \texttt{visualization}. Primary tool records: 58.
\begin{longtable}{p{0.37\linewidth}p{0.47\linewidth}r}
ID & Subcategory & Count\\\hline\endhead
visualization-scientific-3d-visualization-and-postprocessing & Scientific 3D Visualization and Postprocessing & 27\\
visualization-plotting-charts-and-interactive-dashboards & Plotting Charts and Interactive Dashboards & 21\\
visualization-visualization-libraries-and-web-graphics & Visualization Libraries and Web Graphics & 5\\
visualization-specialized-scientific-data-viewers & Specialized Scientific Data Viewers & 13\\
\end{longtable}
\subsection{Rendering and Computer Graphics}
\noindent Domain: \texttt{rendering}. Primary tool records: 93.
\begin{longtable}{p{0.37\linewidth}p{0.47\linewidth}r}
ID & Subcategory & Count\\\hline\endhead
rendering-production & Production and Photorealistic Rendering & 28\\
rendering-architecture & Architectural Visualization & 8\\
rendering-product & Product and Industrial Visualization & 12\\
rendering-realtime & Real-Time Rendering and Interactive Graphics & 20\\
rendering-scientific & Scientific, Spectral and Differentiable Rendering & 19\\
rendering-infrastructure & Rendering Libraries and Scene Infrastructure & 22\\
rendering-content & 3D Scene Creation and Animation & 3\\
rendering-neural & Neural Rendering and Synthetic Scene Generation & 8\\
\end{longtable}
\subsection{Operations Research and Discrete-Event Simulation}
\noindent Domain: \texttt{operations}. Primary tool records: 46.
\begin{longtable}{p{0.37\linewidth}p{0.47\linewidth}r}
ID & Subcategory & Count\\\hline\endhead
operations-discrete-event-and-queueing-simulation & Discrete-Event and Queueing Simulation & 24\\
operations-factory-production-and-material-flow & Factory Production and Material Flow & 7\\
operations-supply-chain-logistics-and-scheduling & Supply Chain Logistics and Scheduling & 13\\
operations-agent-based-and-hybrid-operations-modeling & Agent-Based and Hybrid Operations Modeling & 2\\
operations-hosted-simulation-and-experiment-execution & Hosted Simulation and Experiment Execution & 1\\
operations-factory-layout-and-equipment-planning & Factory Layout and Equipment Planning & 2\\
operations-workstation-ergonomics-and-human-factors & Workstation Ergonomics and Human Factors & 2\\
operations-airport-capacity-and-air-traffic-simulation & Airport Capacity and Air Traffic Simulation & 1\\
operations-engineering-project-portfolio-and-cost-management & Engineering Project Portfolio and Cost Management & 2\\
\end{longtable}
\subsection{Astronomy and Astrophysics}
\noindent Domain: \texttt{astronomy}. Primary tool records: 120.
\begin{longtable}{p{0.37\linewidth}p{0.47\linewidth}r}
ID & Subcategory & Count\\\hline\endhead
astronomy-astronomical-images-and-photometry & Astronomical Images and Photometry & 11\\
astronomy-spectroscopy-and-radio-astronomy & Spectroscopy and Radio Astronomy & 29\\
astronomy-observatory-control-and-observation-planning & Observatory Control and Observation Planning & 17\\
astronomy-astronomical-catalogs-and-data-access & Astronomical Catalogs and Data Access & 39\\
astronomy-cosmology-and-astrophysical-simulation & Cosmology and Astrophysical Simulation & 17\\
astronomy-astronomical-data-reduction-and-statistical-modeling & Astronomical Data Reduction and Statistical Modeling & 11\\
astronomy-astronomical-coordinates-time-and-numerical-utilities & Astronomical Coordinates Time and Numerical Utilities & 1\\
astronomy-astrochemistry-and-radiative-transfer & Astrochemistry and Radiative Transfer & 6\\
astronomy-telescope-instrument-and-observation-simulation & Telescope Instrument and Observation Simulation & 3\\
astronomy-interstellar-chemistry-and-radiative-transfer & Interstellar Chemistry and Radiative Transfer & 1\\
\end{longtable}
\subsection{Experimental Data Analysis and Characterization}
\noindent Domain: \texttt{experimental-analysis}. Primary tool records: 67.
\begin{longtable}{p{0.37\linewidth}p{0.47\linewidth}r}
ID & Subcategory & Count\\\hline\endhead
experimental-analysis-diffraction-crystallography-and-scattering & Diffraction Crystallography and Scattering & 23\\
experimental-analysis-spectroscopy-and-spectral-analysis & Spectroscopy and Spectral Analysis & 18\\
experimental-analysis-microscopy-and-hyperspectral-analysis & Microscopy and Hyperspectral Analysis & 10\\
experimental-analysis-modal-vibration-and-acoustic-testing & Modal Vibration and Acoustic Testing & 3\\
experimental-analysis-experimental-data-reduction-and-fitting & Experimental Data Reduction and Fitting & 3\\
experimental-analysis-optical-flow-measurement-and-particle-velocimetry & Optical Flow Measurement and Particle Velocimetry & 7\\
experimental-analysis-tomographic-reconstruction-and-microscopy-processing & Tomographic Reconstruction and Microscopy Processing & 6\\
experimental-analysis-mass-spectrometry-and-chromatography & Mass Spectrometry and Chromatography & 5\\
experimental-analysis-test-data-management-analysis-and-reporting & Test Data Management Analysis and Reporting & 1\\
\end{longtable}
\subsection{Scientific Imaging and Image Processing}
\noindent Domain: \texttt{imaging}. Primary tool records: 48.
\begin{longtable}{p{0.37\linewidth}p{0.47\linewidth}r}
ID & Subcategory & Count\\\hline\endhead
imaging-image-processing-and-computer-vision & Image Processing and Computer Vision & 20\\
imaging-tomography-and-inverse-reconstruction & Tomography and Inverse Reconstruction & 22\\
imaging-scientific-image-formats-and-acquisition & Scientific Image Formats and Acquisition & 9\\
imaging-3d-photogrammetry-and-image-based-metrology & 3D Photogrammetry and Image-Based Metrology & 1\\
\end{longtable}
\subsection{Scientific Workflow and Computing Infrastructure}
\noindent Domain: \texttt{workflow}. Primary tool records: 56.
\begin{longtable}{p{0.37\linewidth}p{0.47\linewidth}r}
ID & Subcategory & Count\\\hline\endhead
workflow-scientific-workflow-orchestration & Scientific Workflow Orchestration & 34\\
workflow-hpc-scheduling-and-distributed-execution & HPC Scheduling and Distributed Execution & 21\\
workflow-reproducible-research-and-experiment-tracking & Reproducible Research and Experiment Tracking & 7\\
workflow-notebooks-documents-and-scientific-reporting & Notebooks Documents and Scientific Reporting & 9\\
\end{longtable}
\subsection{Geospatial Science and Geographic Information Systems}
\noindent Domain: \texttt{geospatial}. Primary tool records: 53.
\begin{longtable}{p{0.37\linewidth}p{0.47\linewidth}r}
ID & Subcategory & Count\\\hline\endhead
geospatial-desktop-gis-and-spatial-analysis & Desktop GIS and Spatial Analysis & 20\\
geospatial-geospatial-libraries-and-coordinate-systems & Geospatial Libraries and Coordinate Systems & 13\\
geospatial-geospatial-databases-and-web-services & Geospatial Databases and Web Services & 7\\
geospatial-terrain-lidar-and-remote-sensing & Terrain LiDAR and Remote Sensing & 17\\
geospatial-geomagnetic-and-geophysical-reference-computation & Geomagnetic and Geophysical Reference Computation & 1\\
geospatial-geospatial-data-integration-and-etl & Geospatial Data Integration and ETL & 1\\
\end{longtable}
\subsection{Product Lifecycle Management}
\noindent Domain: \texttt{plm}. Primary tool records: 30.
\begin{longtable}{p{0.37\linewidth}p{0.47\linewidth}r}
ID & Subcategory & Count\\\hline\endhead
plm-product-lifecycle-and-engineering-data-management & Product Lifecycle and Engineering Data Management & 22\\
plm-cad-documents-and-design-collaboration & CAD Documents and Design Collaboration & 19\\
plm-bom-configuration-and-product-variability & BOM Configuration and Product Variability & 9\\
\end{longtable}
\subsection{Electrical Systems Design}
\noindent Domain: \texttt{electrical-design}. Primary tool records: 44.
\begin{longtable}{p{0.37\linewidth}p{0.47\linewidth}r}
ID & Subcategory & Count\\\hline\endhead
electrical-design-electrical-schematics-and-panel-engineering & Electrical Schematics and Panel Engineering & 13\\
electrical-design-wire-harness-and-cable-design & Wire Harness and Cable Design & 10\\
electrical-design-electrical-installation-and-cable-sizing & Electrical Installation and Cable Sizing & 9\\
electrical-design-grounding-lightning-and-electromagnetic-compatibility & Grounding Lightning and Electromagnetic Compatibility & 14\\
electrical-design-electromagnetic-field-and-device-analysis & Electromagnetic Field and Device Analysis & 1\\
\end{longtable}
\subsection{Systems Engineering}
\noindent Domain: \texttt{systems-engineering}. Primary tool records: 83.
\begin{longtable}{p{0.37\linewidth}p{0.47\linewidth}r}
ID & Subcategory & Count\\\hline\endhead
systems-engineering-requirements-traceability-and-verification & Requirements Traceability and Verification & 25\\
systems-engineering-mbse-sysml-and-system-architecture & MBSE SysML and System Architecture & 20\\
systems-engineering-system-simulation-and-co-simulation & System Simulation and Co-Simulation & 28\\
systems-engineering-engineering-configuration-and-lifecycle-integration & Engineering Configuration and Lifecycle Integration & 4\\
systems-engineering-formal-requirements-and-architecture-verification & Formal Requirements and Architecture Verification & 4\\
systems-engineering-safety-reliability-and-assurance-cases & Safety Reliability and Assurance Cases & 7\\
systems-engineering-test-automation-and-verification-management & Test Automation and Verification Management & 6\\
systems-engineering-embedded-software-verification-and-assurance & Embedded Software Verification and Assurance & 2\\
systems-engineering-requirements-management-and-traceability & Requirements Management and Traceability & 1\\
\end{longtable}
\subsection{Reliability and Safety Engineering}
\noindent Domain: \texttt{reliability}. Primary tool records: 40.
\begin{longtable}{p{0.37\linewidth}p{0.47\linewidth}r}
ID & Subcategory & Count\\\hline\endhead
reliability-fault-trees-and-system-reliability-modeling & Fault Trees and System Reliability Modeling & 19\\
reliability-failure-modes-safety-and-risk-assessment & Failure Modes Safety and Risk Assessment & 17\\
reliability-life-data-accelerated-testing-and-reliability-prediction & Life Data Accelerated Testing and Reliability Prediction & 8\\
reliability-maintenance-failure-reporting-and-asset-performance & Maintenance Failure Reporting and Asset Performance & 5\\
reliability-risk-uncertainty-and-simulation-frameworks & Risk Uncertainty and Simulation Frameworks & 1\\
\end{longtable}
\subsection{Metrology and Dimensional Inspection}
\noindent Domain: \texttt{metrology}. Primary tool records: 137.
\begin{longtable}{p{0.37\linewidth}p{0.47\linewidth}r}
ID & Subcategory & Count\\\hline\endhead
metrology-coordinate & Coordinate Measurement and Geometric Inspection & 32\\
metrology-optical & Optical and Vision-Based Metrology & 17\\
metrology-scanning & 3D Scanning and Reverse Engineering & 15\\
metrology-surface & Surface Texture and Topography & 12\\
metrology-nde & Nondestructive Testing and Inspection & 45\\
metrology-quality & Measurement Data, Calibration and Quality Analysis & 8\\
metrology-virtual & Virtual Sensing and Model-Based Measurement & 1\\
metrology-dimensional-tolerances-gd-t-and-variation-analysis & Dimensional Tolerances GD\&T and Variation Analysis & 5\\
metrology-calibration-procedures-and-instrument-asset-management & Calibration Procedures and Instrument Asset Management & 9\\
metrology-measurement-uncertainty-and-system-capability & Measurement Uncertainty and System Capability & 8\\
metrology-industrial-machine-vision-and-optical-inspection & Industrial Machine Vision and Optical Inspection & 2\\
metrology-machine-tool-and-coordinate-system-calibration & Machine Tool and Coordinate System Calibration & 3\\
\end{longtable}
\subsection{Industrial Plant and Piping Design}
\noindent Domain: \texttt{plant-design}. Primary tool records: 36.
\begin{longtable}{p{0.37\linewidth}p{0.47\linewidth}r}
ID & Subcategory & Count\\\hline\endhead
plant-design-plant-layout-piping-and-instrumentation-design & Plant Layout Piping and Instrumentation Design & 11\\
plant-design-piping-isometrics-and-fabrication-drawings & Piping Isometrics and Fabrication Drawings & 5\\
plant-design-pipe-stress-and-hydraulic-network-analysis & Pipe Stress and Hydraulic Network Analysis & 9\\
plant-design-pressure-vessel-and-equipment-design & Pressure Vessel and Equipment Design & 11\\
\end{longtable}
\subsection{Infrastructure Operations and Asset Management}
\noindent Domain: \texttt{infrastructure-operations}. Primary tool records: 76.
\begin{longtable}{p{0.37\linewidth}p{0.47\linewidth}r}
ID & Subcategory & Count\\\hline\endhead
infrastructure-operations-civil-infrastructure-and-asset-monitoring & Civil Infrastructure and Asset Monitoring & 26\\
infrastructure-operations-transportation-and-traffic-simulation & Transportation and Traffic Simulation & 23\\
infrastructure-operations-water-wastewater-and-stormwater-operations & Water Wastewater and Stormwater Operations & 12\\
infrastructure-operations-site-design-grading-and-earthworks & Site Design Grading and Earthworks & 1\\
infrastructure-operations-hydraulic-infrastructure-design & Hydraulic Infrastructure Design & 3\\
infrastructure-operations-infrastructure-risk-and-incident-modeling & Infrastructure Risk and Incident Modeling & 1\\
infrastructure-operations-road-rail-and-corridor-design & Road Rail and Corridor Design & 7\\
infrastructure-operations-pavement-condition-and-maintenance-planning & Pavement Condition and Maintenance Planning & 2\\
infrastructure-operations-construction-machine-control-and-survey-guidance & Construction Machine Control and Survey Guidance & 1\\
infrastructure-operations-bridge-condition-monitoring-and-asset-management & Bridge Condition Monitoring and Asset Management & 2\\
infrastructure-operations-transportation-asset-data-and-linear-referencing & Transportation Asset Data and Linear Referencing & 1\\
infrastructure-operations-flood-consequences-and-evacuation-simulation & Flood Consequences and Evacuation Simulation & 1\\
\end{longtable}
\subsection{Manufacturing Cost Estimation}
\noindent Domain: \texttt{manufacturing-cost}. Primary tool records: 27.
\begin{longtable}{p{0.37\linewidth}p{0.47\linewidth}r}
ID & Subcategory & Count\\\hline\endhead
manufacturing-cost-manufacturing-cost-estimation-and-quoting & Manufacturing Cost Estimation and Quoting & 20\\
manufacturing-cost-design-for-manufacturing-and-process-assessment & Design for Manufacturing and Process Assessment & 8\\
\end{longtable}
\subsection{Cloud Engineering and Simulation Platforms}
\noindent Domain: \texttt{cloud-platform}. Primary tool records: 35.
\begin{longtable}{p{0.37\linewidth}p{0.47\linewidth}r}
ID & Subcategory & Count\\\hline\endhead
cloud-platform-hpc-provisioning-and-engineering-job-execution & HPC Provisioning and Engineering Job Execution & 26\\
cloud-platform-hosted-engineering-applications-and-apis & Hosted Engineering Applications and APIs & 6\\
cloud-platform-remote-visualization-and-engineering-desktops & Remote Visualization and Engineering Desktops & 4\\
cloud-platform-scientific-software-builds-and-environment-management & Scientific Software Builds and Environment Management & 2\\
\end{longtable}
\subsection{Engineering Collaboration and Design Review}
\noindent Domain: \texttt{collaboration}. Primary tool records: 20.
\begin{longtable}{p{0.37\linewidth}p{0.47\linewidth}r}
ID & Subcategory & Count\\\hline\endhead
collaboration-engineering-model-review-and-visualization & Engineering Model Review and Visualization & 11\\
collaboration-engineering-information-and-project-collaboration & Engineering Information and Project Collaboration & 9\\
collaboration-technical-publishing-and-interactive-applications & Technical Publishing and Interactive Applications & 5\\
\end{longtable}
\subsection{Plasma Physics and Engineering}
\noindent Domain: \texttt{plasma}. Primary tool records: 67.
\begin{longtable}{p{0.37\linewidth}p{0.47\linewidth}r}
ID & Subcategory & Count\\\hline\endhead
plasma-pic & Particle-in-Cell and Kinetic Plasma Simulation & 23\\
plasma-gyrokinetic & Gyrokinetic Turbulence and Transport & 8\\
plasma-mhd & Magnetohydrodynamics and Plasma Stability & 8\\
plasma-equilibrium & Magnetic Equilibrium and Confinement Design & 12\\
plasma-edge & Plasma Edge, Divertors and Neutral Transport & 3\\
plasma-discharge & Low-Temperature Plasmas and Industrial Discharges & 11\\
plasma-integrated & Integrated Fusion and Core Transport Modeling & 1\\
plasma-fluid & Plasma Fluid Modeling Frameworks & 1\\
plasma-acceleration & Plasma Acceleration and Wakefields & 2\\
plasma-electric-propulsion-and-plasma-thrusters & Electric Propulsion and Plasma Thrusters & 1\\
plasma-energetic-particle-orbits-and-neutral-beam-injection & Energetic Particle Orbits and Neutral-Beam Injection & 2\\
plasma-disruptions-runaway-electrons-and-plasma-startup & Disruptions Runaway Electrons and Plasma Startup & 2\\
plasma-atomic-kinetics-and-plasma-spectroscopy & Atomic Kinetics and Plasma Spectroscopy & 3\\
plasma-laser-pulse-preparation-and-simulation-interfaces & Laser Pulse Preparation and Simulation Interfaces & 1\\
\end{longtable}
\subsection{Quantum Many-Body and Open-System Simulation}
\noindent Domain: \texttt{quantum-manybody}. Primary tool records: 49.
\begin{longtable}{p{0.37\linewidth}p{0.47\linewidth}r}
ID & Subcategory & Count\\\hline\endhead
quantum-manybody-tensor-networks-and-dmrg & Tensor Networks and DMRG & 19\\
quantum-manybody-quantum-monte-carlo-and-variational-states & Quantum Monte Carlo and Variational States & 7\\
quantum-manybody-exact-diagonalization-and-lattice-models & Exact Diagonalization and Lattice Models & 12\\
quantum-manybody-open-quantum-systems-and-quantum-dynamics & Open Quantum Systems and Quantum Dynamics & 5\\
quantum-manybody-strong-correlations-and-dynamical-mean-field-theory & Strong Correlations and Dynamical Mean-Field Theory & 8\\
quantum-manybody-superconductivity-and-mesoscopic-quantum-systems & Superconductivity and Mesoscopic Quantum Systems & 1\\
quantum-manybody-many-body-electronic-excitations & Many-Body Electronic Excitations & 1\\
\end{longtable}
\subsection{Quantum Computing and Circuit Simulation}
\noindent Domain: \texttt{quantum-computing}. Primary tool records: 42.
\begin{longtable}{p{0.37\linewidth}p{0.47\linewidth}r}
ID & Subcategory & Count\\\hline\endhead
quantum-computing-quantum-circuit-simulation-and-sdks & Quantum Circuit Simulation and SDKs & 27\\
quantum-computing-quantum-error-correction-and-stabilizer-simulation & Quantum Error Correction and Stabilizer Simulation & 4\\
quantum-computing-photonic-and-continuous-variable-quantum-computing & Photonic and Continuous-Variable Quantum Computing & 3\\
quantum-computing-quantum-algorithms-and-hybrid-optimization & Quantum Algorithms and Hybrid Optimization & 3\\
quantum-computing-quantum-hardware-control-and-pulse-modeling & Quantum Hardware Control and Pulse Modeling & 2\\
quantum-computing-quantum-networks-and-communication & Quantum Networks and Communication & 4\\
quantum-computing-annealing-and-binary-optimization & Annealing and Binary Optimization & 2\\
quantum-computing-quantum-processor-layout-and-design & Quantum Processor Layout and Design & 2\\
\end{longtable}
\subsection{Particle Accelerator and Beam Physics}
\noindent Domain: \texttt{accelerator}. Primary tool records: 40.
\begin{longtable}{p{0.37\linewidth}p{0.47\linewidth}r}
ID & Subcategory & Count\\\hline\endhead
accelerator-beam-optics-and-lattice-design & Beam Optics and Lattice Design & 13\\
accelerator-particle-tracking-and-beam-dynamics & Particle Tracking and Beam Dynamics & 29\\
accelerator-accelerator-electromagnetics-and-wakefields & Accelerator Electromagnetics and Wakefields & 7\\
accelerator-charged-particle-optics-and-sources & Charged-Particle Optics and Sources & 4\\
accelerator-synchrotron-radiation-and-beamline-engineering & Synchrotron Radiation and Beamline Engineering & 1\\
\end{longtable}
\subsection{Fire and Process Safety Engineering}
\noindent Domain: \texttt{fire-safety}. Primary tool records: 27.
\begin{longtable}{p{0.37\linewidth}p{0.47\linewidth}r}
ID & Subcategory & Count\\\hline\endhead
fire-safety-chemical-release-consequences-and-quantitative-risk & Chemical Release Consequences and Quantitative Risk & 6\\
fire-safety-wildland-fire-behavior-and-atmosphere-coupling & Wildland Fire Behavior and Atmosphere Coupling & 6\\
fire-safety-fire-protection-design-and-code-compliance & Fire Protection Design and Code Compliance & 1\\
fire-safety-fire-smoke-and-ventilation-simulation & Fire Smoke and Ventilation Simulation & 2\\
fire-safety-pedestrian-movement-and-evacuation & Pedestrian Movement and Evacuation & 4\\
fire-safety-structural-response-to-fire & Structural Response to Fire & 4\\
fire-safety-sprinkler-design-and-hydraulic-calculation & Sprinkler Design and Hydraulic Calculation & 3\\
fire-safety-smoke-control-and-building-airflow & Smoke Control and Building Airflow & 1\\
\end{longtable}
\subsection{Planetary Science}
\noindent Domain: \texttt{planetary-science}. Primary tool records: 14.
\begin{longtable}{p{0.37\linewidth}p{0.47\linewidth}r}
ID & Subcategory & Count\\\hline\endhead
planetary-science-planetary-atmosphere-spectra-and-retrieval & Planetary Atmosphere Spectra and Retrieval & 4\\
planetary-science-surface-chronology-and-crater-analysis & Surface Chronology and Crater Analysis & 1\\
planetary-science-planetary-climate-and-general-circulation & Planetary Climate and General Circulation & 3\\
planetary-science-planetary-imaging-mapping-and-mission-data & Planetary Imaging Mapping and Mission Data & 2\\
planetary-science-engineering-atmosphere-and-environment-models & Engineering Atmosphere and Environment Models & 2\\
planetary-science-planetary-surface-thermal-modeling & Planetary Surface Thermal Modeling & 1\\
planetary-science-impact-cratering-and-shock-physics & Impact Cratering and Shock Physics & 1\\
\end{longtable}
\subsection{Solid Earth Sciences}
\noindent Domain: \texttt{geoscience}. Primary tool records: 58.
\begin{longtable}{p{0.37\linewidth}p{0.47\linewidth}r}
ID & Subcategory & Count\\\hline\endhead
geoscience-landscape-evolution-and-sedimentary-basins & Landscape Evolution and Sedimentary Basins & 4\\
geoscience-mineral-thermodynamics-and-petrology & Mineral Thermodynamics and Petrology & 5\\
geoscience-crustal-deformation-and-earthquake-mechanics & Crustal Deformation and Earthquake Mechanics & 2\\
geoscience-geological-modeling-and-subsurface-interpretation & Geological Modeling and Subsurface Interpretation & 6\\
geoscience-geochronology-and-paleoclimate-age-models & Geochronology and Paleoclimate Age Models & 5\\
geoscience-gnss-geodesy-and-tectonic-positioning & GNSS Geodesy and Tectonic Positioning & 4\\
geoscience-mantle-convection-and-lithosphere-dynamics & Mantle Convection and Lithosphere Dynamics & 5\\
geoscience-structural-geology-and-orientation-analysis & Structural Geology and Orientation Analysis & 3\\
geoscience-plate-reconstruction-and-tectonic-evolution & Plate Reconstruction and Tectonic Evolution & 3\\
geoscience-geophysical-inversion-and-potential-fields & Geophysical Inversion and Potential Fields & 4\\
geoscience-seismic-data-processing-and-monitoring & Seismic Data Processing and Monitoring & 6\\
geoscience-earthquake-location-hazard-and-risk & Earthquake Location Hazard and Risk & 3\\
geoscience-seismic-wave-propagation-and-dynamic-rupture & Seismic Wave Propagation and Dynamic Rupture & 3\\
geoscience-geodynamics-modeling-frameworks & Geodynamics Modeling Frameworks & 1\\
geoscience-geostatistics-and-spatial-property-modeling & Geostatistics and Spatial Property Modeling & 3\\
geoscience-geochronology-and-geochemical-data-analysis & Geochronology and Geochemical Data Analysis & 1\\
\end{longtable}
\subsection{Microfluidic Engineering}
\noindent Domain: \texttt{microfluidics}. Primary tool records: 7.
\begin{longtable}{p{0.37\linewidth}p{0.47\linewidth}r}
ID & Subcategory & Count\\\hline\endhead
microfluidics-microfluidic-device-design-and-automation & Microfluidic Device Design and Automation & 5\\
microfluidics-microscale-flow-transport-and-electrokinetics & Microscale Flow Transport and Electrokinetics & 2\\
\end{longtable}
\subsection{Data Center Engineering}
\noindent Domain: \texttt{data-center}. Primary tool records: 12.
\begin{longtable}{p{0.37\linewidth}p{0.47\linewidth}r}
ID & Subcategory & Count\\\hline\endhead
data-center-cooling-airflow-and-thermal-design & Cooling Airflow and Thermal Design & 4\\
data-center-infrastructure-capacity-layout-and-asset-planning & Infrastructure Capacity Layout and Asset Planning & 6\\
data-center-operational-thermal-and-power-optimization & Operational Thermal and Power Optimization & 1\\
data-center-workload-scheduling-and-energy-simulation & Workload Scheduling and Energy Simulation & 1\\
\end{longtable}
\subsection{Agricultural Engineering}
\noindent Domain: \texttt{agricultural}. Primary tool records: 5.
\begin{longtable}{p{0.37\linewidth}p{0.47\linewidth}r}
ID & Subcategory & Count\\\hline\endhead
agricultural-greenhouse-climate-and-energy-engineering & Greenhouse Climate and Energy Engineering & 1\\
agricultural-agricultural-application-and-spray-engineering & Agricultural Application and Spray Engineering & 1\\
agricultural-irrigation-design-and-hydraulics & Irrigation Design and Hydraulics & 3\\
\end{longtable}
\subsection{Naval Architecture and Marine Engineering}
\noindent Domain: \texttt{marine}. Primary tool records: 73.
\begin{longtable}{p{0.37\linewidth}p{0.47\linewidth}r}
ID & Subcategory & Count\\\hline\endhead
marine-seakeeping-wave-loads-and-potential-flow-hydrodynamics & Seakeeping Wave Loads and Potential-Flow Hydrodynamics & 16\\
marine-mooring-risers-cables-and-offshore-dynamics & Mooring Risers Cables and Offshore Dynamics & 6\\
marine-marine-operations-and-coupled-vessel-dynamics & Marine Operations and Coupled Vessel Dynamics & 9\\
marine-vessel-resistance-propellers-and-propulsion & Vessel Resistance Propellers and Propulsion & 12\\
marine-hydrodynamic-model-preparation-and-data-exchange & Hydrodynamic Model Preparation and Data Exchange & 2\\
marine-hull-structural-design-and-classification-rules & Hull Structural Design and Classification Rules & 8\\
marine-integrated-ship-design-and-production-engineering & Integrated Ship Design and Production Engineering & 8\\
marine-hull-form-design-and-fairing & Hull Form Design and Fairing & 7\\
marine-hydrostatics-stability-and-loading & Hydrostatics Stability and Loading & 11\\
marine-viscous-flow-and-marine-cfd & Viscous Flow and Marine CFD & 4\\
\end{longtable}
\subsection{Mining and Mineral Processing}
\noindent Domain: \texttt{mining}. Primary tool records: 32.
\begin{longtable}{p{0.37\linewidth}p{0.47\linewidth}r}
ID & Subcategory & Count\\\hline\endhead
mining-drill-and-blast-design-and-analysis & Drill-and-Blast Design and Analysis & 6\\
mining-mineral-processing-and-metallurgical-flowsheets & Mineral Processing and Metallurgical Flowsheets & 9\\
mining-mine-design-production-planning-and-scheduling & Mine Design Production Planning and Scheduling & 8\\
mining-geological-models-and-mineral-resource-estimation & Geological Models and Mineral Resource Estimation & 4\\
mining-haulage-fleet-and-mine-logistics-simulation & Haulage Fleet and Mine Logistics Simulation & 1\\
mining-exploration-drillhole-and-assay-data-management & Exploration Drillhole and Assay Data Management & 1\\
mining-mine-ventilation-cooling-and-environmental-control & Mine Ventilation Cooling and Environmental Control & 3\\
\end{longtable}
\subsection{Urban and Regional Planning}
\noindent Domain: \texttt{urban-planning}. Primary tool records: 16.
\begin{longtable}{p{0.37\linewidth}p{0.47\linewidth}r}
ID & Subcategory & Count\\\hline\endhead
urban-planning-urban-massing-zoning-and-site-feasibility & Urban Massing Zoning and Site Feasibility & 4\\
urban-planning-semantic-city-models-and-geospatial-integration & Semantic City Models and Geospatial Integration & 2\\
urban-planning-land-use-scenarios-and-spatial-suitability & Land-Use Scenarios and Spatial Suitability & 1\\
urban-planning-urban-microclimate-and-environmental-simulation & Urban Microclimate and Environmental Simulation & 6\\
urban-planning-urban-energy-and-infrastructure-planning & Urban Energy and Infrastructure Planning & 1\\
urban-planning-urban-land-use-and-development-microsimulation & Urban Land-Use and Development Microsimulation & 2\\
\end{longtable}
\subsection{Mechanical Systems and Machine Elements}
\noindent Domain: \texttt{mechanical-systems}. Primary tool records: 70.
\begin{longtable}{p{0.37\linewidth}p{0.47\linewidth}r}
ID & Subcategory & Count\\\hline\endhead
mechanical-systems-bearings-lubrication-and-rolling-contacts & Bearings Lubrication and Rolling Contacts & 16\\
mechanical-systems-powertrain-dynamics-durability-and-nvh & Powertrain Dynamics Durability and NVH & 7\\
mechanical-systems-contact-mechanics-friction-and-wear & Contact Mechanics Friction and Wear & 5\\
mechanical-systems-gears-transmissions-and-machine-element-design & Gears Transmissions and Machine-Element Design & 33\\
mechanical-systems-mechanical-seal-simulation & Mechanical Seal Simulation & 2\\
mechanical-systems-turbomachinery-meanline-and-throughflow-design & Turbomachinery Meanline and Throughflow Design & 3\\
mechanical-systems-shafts-bearings-and-drivetrain-analysis & Shafts Bearings and Drivetrain Analysis & 6\\
mechanical-systems-gear-metrology-and-manufacturing-verification & Gear Metrology and Manufacturing Verification & 2\\
\end{longtable}
\subsection{Atmospheric Science and Meteorology}
\noindent Domain: \texttt{atmospheric-science}. Primary tool records: 65.
\begin{longtable}{p{0.37\linewidth}p{0.47\linewidth}r}
ID & Subcategory & Count\\\hline\endhead
water-atmosphere & Atmospheric Composition and Air Quality & 26\\
water-weather & Weather Prediction and Climate Modeling & 21\\
atmospheric-science-weather-and-climate-data-processing & Weather and Climate Data Processing & 4\\
atmospheric-science-atmospheric-radiative-transfer & Atmospheric Radiative Transfer & 7\\
atmospheric-science-earth-system-model-coupling-and-infrastructure & Earth-System Model Coupling and Infrastructure & 1\\
atmospheric-science-earth-system-data-assimilation & Earth-System Data Assimilation & 4\\
atmospheric-science-forecast-verification-and-evaluation & Forecast Verification and Evaluation & 2\\
\end{longtable}
\subsection{Embedded Software Engineering}
\noindent Domain: \texttt{embedded-software}. Primary tool records: 68.
\begin{longtable}{p{0.37\linewidth}p{0.47\linewidth}r}
ID & Subcategory & Count\\\hline\endhead
embedded-software-embedded-ides-and-development-toolchains & Embedded IDEs and Development Toolchains & 12\\
embedded-software-compilers-and-cross-compilation & Compilers and Cross Compilation & 4\\
embedded-software-real-time-operating-systems-and-partitioning & Real-Time Operating Systems and Partitioning & 12\\
embedded-software-autosar-architecture-and-basic-software & AUTOSAR Architecture and Basic Software & 7\\
embedded-software-virtual-ecus-and-embedded-system-simulation & Virtual ECUs and Embedded System Simulation & 7\\
embedded-software-ecu-measurement-calibration-and-diagnostics & ECU Measurement Calibration and Diagnostics & 2\\
embedded-software-model-based-embedded-code-generation & Model-Based Embedded Code Generation & 2\\
embedded-software-static-analysis-and-software-verification & Static Analysis and Software Verification & 9\\
embedded-software-embedded-debugging-trace-and-profiling & Embedded Debugging Trace and Profiling & 4\\
embedded-software-embedded-linux-distribution-and-image-builds & Embedded Linux Distribution and Image Builds & 2\\
embedded-software-embedded-unit-integration-and-regression-testing & Embedded Unit Integration and Regression Testing & 2\\
embedded-software-build-systems-and-dependency-management & Build Systems and Dependency Management & 5\\
\end{longtable}

\end{document}
